% Generated draft: requires mathematical and source review.
\documentclass[11pt]{article}
\usepackage[margin=25mm]{geometry}
\usepackage{fontspec}
\setmainfont{Latin Modern Roman}
\usepackage{amsmath,amssymb,mathtools,unicode-math}
\setmathfont{Latin Modern Math}
\usepackage{enumitem,needspace,xurl,hyperref}
\setlength{\parindent}{0pt}
\setlength{\parskip}{5pt}
\setlength{\emergencystretch}{3em}
\newcommand{\N}{\mathbb{N}}
\newcommand{\Z}{\mathbb{Z}}
\newcommand{\Q}{\mathbb{Q}}
\newcommand{\R}{\mathbb{R}}
\newcommand{\C}{\mathbb{C}}
\newcommand{\F}{\mathbb{F}}
\newcommand{\A}{\mathbb{A}}
\newcommand{\PP}{\mathbb{P}}
\newcommand{\E}{\mathbb{E}}
\newcommand{\eps}{\varepsilon}
\newcommand{\dd}{\,\mathrm{d}}
\begin{document}
\section*{UnsolvedMath ID: 701612}
\textbf{Source problem code:} catalogue entry
\subsection{Definitions and mathematical statement}
Work in the first-order language of groups \(\mathcal L_{\mathrm{grp}}=\{\cdot,{}^{-1},1\}\), with no additional parameters. For each positive integer \(n\), let \(\operatorname{Form}_n\) be the set of formulas in this language whose free variables are among \(x_1,\ldots,x_n\). For a group \(G\) and a tuple \(\bar g=(g_1,\ldots,g_n)\in G^n\), define
\[
\operatorname{Tp}^{G}(\bar g)=\{\varphi\in\operatorname{Form}_n:G\models\varphi(\bar g)\}.
\]
Groups \(G\) and \(H\) are \emph{isotypic} if, for every integer \(n\geq 1\),
\[
\begin{aligned}
&\forall\bar g\in G^n\;\exists\bar h\in H^n:
\operatorname{Tp}^{G}(\bar g)=\operatorname{Tp}^{H}(\bar h),\\
&\forall\bar h\in H^n\;\exists\bar g\in G^n:
\operatorname{Tp}^{H}(\bar h)=\operatorname{Tp}^{G}(\bar g).
\end{aligned}
\]
A group is \emph{finitely generated} if it is generated by a finite subset.

\textbf{Question (B. I. Plotkin, Kourovka Notebook Problem 19.64).} For all finitely generated groups \(G\) and \(H\), does isotypy imply that \(G\cong H\)? Equivalently, under these hypotheses, must there exist a bijection \(f:G\to H\) satisfying \(f(ab)=f(a)f(b)\) for all \(a,b\in G\)?

\textbf{Supplied contextual claims, requiring review.} The record labels the question open and reports a positive answer when one of the finitely generated groups is free (R. Sklinos), abelian (G. Zhitomirski), virtually polycyclic, metabelian, free solvable (A. Myasnikov and N. Romanovskii), co-Hopfian (with the parenthetical assertion ``in particular homogeneous''), finitely presented Hopfian, or geometrically Noetherian Hopfian (R. Sklinos). The placement and scope of these attributions are retained from the supplied text and have not been independently checked.

The record further attributes to R. Sklinos the assertion that torsion-free hyperbolic groups, braid groups, linear groups, mapping class groups of compact surfaces, limit groups, quasi-cyclic groups, \(\operatorname{Out}(F_n)\) for \(n>2\), and irreducible arithmetic lattices that are not virtually free in semi-simple Lie groups are ``of such kind.'' Here \(F_n\) denotes the free group of rank \(n\), and \(\operatorname{Out}(F_n)\) its automorphism group modulo inner automorphisms. The phrase ``of such kind'' does not specify which preceding condition applies to each class. The inclusion of quasi-cyclic groups, which are not finitely generated, requires clarification in relation to the finite-generation hypothesis. The meaning of ``homogeneous'' and its asserted relationship to the preceding condition also require clarification. These contextual statements are preserved as supplied claims, without adding hypotheses or treating them as verified results.
\subsection{Short English statement}
Must any two finitely generated groups be isomorphic if they realize exactly the same complete first-order types of finite tuples in the language of groups without parameters? This question is attributed to B. I. Plotkin. The supplied record labels it open; its reported positive cases and examples require source and scope review.
\subsection{Sources}
\begin{itemize}
\item Local UnsolvedMath registry, supplied complete record: permanent ID \(701612\), exact source code \texttt{catalogue entry}, title \emph{Kourovka Notebook Problem 19.64}. The record attributes the problem to B. I. Plotkin and gives the proposed year as 2018. Its supplied status is \texttt{open}.
\item Supplied source attribution: \emph{The Kourovka Notebook (Unsolved Problems in Group Theory)}, 19th issue (2018), Problem 19.64; also identified in the record with the 21st issue edition (2026), \texttt{arXiv:1401.0300v48}. The collection description names E. I. Khukhro and V. D. Mazurov as editors. No problem-specific page, theorem citation, or external URL is supplied.
\item The supplied statement attributes reported positive cases to R. Sklinos, G. Zhitomirski, A. Myasnikov, and N. Romanovskii as recorded above. It supplies no separate bibliographic references for those results.
\item Registry update timestamp: \texttt{2026-10-16T14:00:00Z}; upstream revision: \texttt{d371ee917ac9647fbb06cd6fc529deec6f3dda82}. This is supplied registry metadata, not an independently verified literature assessment. No \texttt{literature\_checked\_at} date or \texttt{literature\_sources} links are supplied. No external source was consulted for this draft.
\end{itemize}
\end{document}
